\section{Experiments}
\label{sec:exp}

\subsection{Setup}
\label{sec:exp-setup}
\paragraph{Data.}
\textbf{Open LLM Leaderboard (OLLB)}: the public \texttt{contents} snapshot (4576 rows; 4497 models after removing duplicate names, 4496 after excluding one flagged model) with six tasks (IFEval, BBH, MATH Lvl~5, GPQA, MuSR, MMLU-Pro) in the leaderboard's own normalised scale.
\textbf{HELM MMLU}: release v1.13.0, 79 models on 57 subjects (accuracy). Versions, commits and checksums are in Appendix~\ref{app:repro}.
Because OLLB scores are rescaled per task (e.g.\ GPQA tops out near 29), scale-sensitive rules depend on normalisation, so we also run the raw scores (Section~\ref{sec:exp-robust}).

\paragraph{Targets, rules, caps.}
Every model of rank $\ge 2$ under the mean is a target: 4495 on OLLB and 78 on HELM.
Score rules (mean, geometric, harmonic, median, min) run on all targets. Borda, Condorcet-winner, plurality and the trimmed mean are exact but expensive on OLLB, so they run on a rank-stratified sample of 57 OLLB targets (ranks 2--21 plus a log-uniform sample of the rest, seed 0); on HELM every rule runs on every target.
Cross-rule comparisons therefore use either the 20 closest challengers (ranks 2--21), on which every rule was run, or ratios paired per target.
A target of mean-rule rank $\ge2$ can already be first under another rule (at most one per rule and dataset: HELM geometric, harmonic, min and plurality; OLLB median and min; one Borda target in the sample); its $K$ and $L_1$ are 0 and it is \emph{included} in Table~\ref{tab:main}, Tables~\ref{tab:A1ollb}--\ref{tab:A2helm}, Figure~\ref{fig:rank} and the slopes, but excluded from the uniform-plan ratios (Figure~\ref{fig:uniform}).
We present caps $c_t$ as defined in Section~\ref{sec:method}; the headline cap is the empirical one (\texttt{emp95}), the 95th percentile of the gain a fine-tuned model shows over its base model on that task, measured on 540 base/fine-tune pairs in OLLB, and \texttt{ceil} is the permissive bound.
HELM has no base/fine-tune pairs, so its \texttt{emp95} cap is the OLLB MMLU-Pro quantile (about 2.9 raw points) transferred as a uniform per-subject cap; every HELM statement under that cap is conditional on this transfer.
Margin $\varepsilon=10^{-3}$ points.

\paragraph{Uncertainty.}
Group medians carry 95\% percentile bootstrap intervals over targets (1000 resamples, seed 0); feasibility fractions carry Wilson intervals; medians of $K$ and $L_1$ are over feasible targets and say so.
Targets in neighbouring ranks are not independent (many OLLB entries are merges or fine-tunes of the same base), so intervals describe sampling of targets from this leaderboard, not of leaderboards; we add leaderboard slices below.

\subsection{Theory anchor: closed forms agree with the solvers}
\label{sec:exp-theory}
For the mean, $L_1^\ast=n(\tau+\varepsilon)-\sum_t x_{it}$ matches the solver on every feasible row to $10^{-11}$ points or better; for median and min the closed-form predictions for $K$ (and $L_1$ where derived) match on every row where a prediction exists (Figure~\ref{fig:theory} in Appendix~\ref{app:more}: 15\,666 mean, 4865 median and 17\,550 min rows for $K$, all equal).
Two of our prior expectations are narrowed or contradicted by the data.
(i) The median needs at least $\lfloor n/2\rfloor+1$ tasks \emph{only in the worst case}, when the target is below the rivals' decisive score on every task. Among the 20 closest challengers under the ceiling the median needs $K=1$ on OLLB ($n=6$) and 10 [6.5, 10.5] on HELM ($n=57$), the min needs 2 and 3.5 [2, 5]: these challengers already beat the decisive scores on most tasks.
(ii) The median and the min do not cost as many points as the mean, but \emph{fewer}: the paired ratio $L_1(\text{median})/L_1(\text{mean})$ is 0.60 [0.60, 0.60] on OLLB and 0.156 [0.14, 0.19] on HELM, and for min 0.251 [0.248, 0.253] and 0.147 [0.12, 0.18] (ceiling, all targets).
The mean must buy the whole gap to the leader times $n$ in points; order-statistic rules only fix the few decisive tasks.
``Raise only one task'' is likewise a property of the head of a 6-task leaderboard: under the ceiling the mean needs $K=1$ for the 20 closest OLLB challengers, but only 9.2\% of all OLLB targets, and $K=7$ [4.5, 9] on HELM.

\subsection{Rule ordering of the minimum investment}
\label{sec:exp-rules}

\paragraph{Under the permissive cap the mean is the most expensive rule in points.}
For the 20 closest OLLB challengers (ranks 2--21, ceiling), the median $L_1$ is: median rule 3.6 [1.3, 9.0], Borda 6.3 [4.8, 8.8], min 10.4 [6.3, 15.2], Condorcet winner 11.4 [10.4, 14.4], plurality 16.4 [14.7, 19.1], harmonic 18.0 [13.3, 20.8], geometric 21.9 [18.8, 23.8], trimmed mean 24.4 [21.3, 26.0], mean 29.2 [26.4, 31.4].
On HELM the order is plurality 16.0 [6.5, 21.1], min 17.0 [6.9, 25.3] and median 17.8 [11.4, 34.0] (intervals overlap), then Condorcet 64.8 [24.5, 94.1], Borda 147 [55, 206], harmonic 242, geometric 258 and mean 278 [177, 360].
Paired on the same targets, the ratio of each rule's $L_1$ to the mean's (ceiling, all targets) is, on HELM, 0.961 [0.953, 0.967] (geometric), 0.925 (harmonic), 0.776 [0.725, 0.823] (Borda), 0.295 [0.282, 0.307] (Condorcet), 0.156 (median), 0.147 (min) and 0.068 [0.063, 0.073] (plurality); on OLLB 0.852 (geometric), 0.697 (harmonic), 0.599 (median), 0.251 (min), and for the 57-target sample 0.372 [0.333, 0.436] (Borda), 0.640 [0.560, 0.655] (Condorcet), 0.500 [0.447, 0.572] (plurality) and 0.815 (trimmed mean).
Table~\ref{tab:main} gives the headline-cap version and Figure~\ref{fig:rank} shows how $L_1$ and feasibility depend on the target's rank.

\begin{table}[t]
\caption{Minimum investment under the headline cap (\texttt{emp95}) for the 20 closest challengers (ranks 2--21) and all other targets: median tasks touched $K$ and points $L_1$ [95\% bootstrap CI] over feasible targets, and number of infeasible targets out of all (a target already first under the rule counts with $K=L_1=0$). HELM uses the cap transferred from OLLB; $^\dagger$ rules run on the 57-target OLLB sample (trimmed mean has no $K$); single numbers have fewer than five feasible targets.}
\label{tab:main}
\centering\small
\input{../tables/table1_rule_dataset.tex}
\end{table}

\paragraph{Where this ordering fails or is unresolved.}
(i) For rank 2--5 there are only four targets per dataset; intervals straddle 1 (for example OLLB ceiling Condorcet 0.91 [0.41, 1.41], plurality 0.84 [0.59, 1.09]), so we make no ordering claim there.
(ii) Under the empirical cap the geometric mean is not cheaper than the mean where both are feasible (1.05 [1.00, 1.07], 111 of 4495 targets), and on HELM at a fixed $+10$ cap geometric and harmonic means are indistinguishable from it (0.99 [0.97, 1.04], 0.999 [0.93, 1.03]).
(iii) \emph{Breadth is a different statement from volume}: on HELM, Borda needs 28.5 [27, 29] \emph{more} tasks than the mean, Condorcet 12 [9, 14] more, median 4 [3, 4] more, while min needs 4 [3, 4] fewer and plurality 7 [4, 9.5] fewer; on OLLB the differences are 0--2 tasks. Borda and Condorcet are cheaper in points but spread the raise over many of the 57 subjects.
(iv) The ordering depends on normalisation and on the task set (Section~\ref{sec:exp-robust}).

\paragraph{Investment grows with rank.}
Under the ceiling, the log-log slope of $L_1$ against the target's mean-rule rank (feasible targets of rank $\ge 2$, bootstrap CI) is 0.44 [0.43, 0.44] (mean) on OLLB, with 0.51 [0.50, 0.52] (geometric), 0.64 [0.62, 0.67] (median), 0.61 [0.58, 0.63] (min), 0.60 [0.54, 0.66] (Borda), 0.44 [0.42, 0.46] (Condorcet) and 0.24 [0.20, 0.29] (plurality); on HELM it is 1.13 [0.91, 1.32] (mean), 1.07 (geometric), 2.17 [1.63, 2.76] (median), 1.76 [1.17, 2.41] (Condorcet), 1.67 [1.48, 1.89] (min), 1.48 [1.25, 1.80] (Borda) and 1.52 [1.12, 2.08] (plurality).
Under capped scenarios the slopes are fitted only to the feasible survivors, which biases them towards flatter values, so we do not interpret them as elasticities.


\begin{figure}[t]
\centering
\includegraphics[width=\linewidth]{../figures/fig1_investment_vs_rank_compact.pdf}
\caption{Investment against the target's current rank under the mean rule, headline cap (OLLB \texttt{emp95}; HELM \texttt{emp95} transferred from OLLB$^\dagger$). Top: fraction of targets for which first place is feasible. Bottom: median score points $L_1$ over feasible targets per rank bin (the median $K$ is in Table~\ref{tab:main}) (bands: 95\% bootstrap CI where at least five targets are feasible). Dashed lines: rules run on the 57-target OLLB sample. HELM bins hold 4, 15 and 59 targets. $^\dagger$HELM has no base/fine-tune pairs.}
\label{fig:rank}
\end{figure}

\subsection{Caps decide who can reach first place at all}
\label{sec:exp-caps}
Under the ceiling, front ($+0$, $+2$) caps no target on either leaderboard is infeasible for any rule.
Under the empirical cap the situation reverses (Table~\ref{tab:feas}, Figure~\ref{fig:caps}): on OLLB, first place under the mean is infeasible for 92.5\% of targets (nearly all with rank $>100$; the smallest infeasible rank is 110), whereas min, Condorcet and harmonic have infeasible targets already at ranks 5--6 and median at 13.
\emph{The ordering of rules by cost flips with the cap}: the mean is the most expensive rule in points under the ceiling, but under \texttt{emp95} it is the most feasible rule beyond rank 20, because the compensatory mean can spread a small gain over six tasks while min, median and the Condorcet requirement need specific tasks lifted past specific rivals.
This is dataset-dependent: on HELM under a $+10$ cap plurality (27\% infeasible) and median (49\%) are the most feasible rules and min the least (82\%); under the transferred \texttt{emp95} 85--96\% of targets are infeasible for every rule, and the smallest infeasible rank is 3 (min) to 10 (median).
So ``any model can reach first place'' (ceiling) becomes ``only models within roughly the top 100 (OLLB, mean) or top 6 (HELM, mean) can'' under the empirical caps.
The fractions for Borda, trimmed mean, plurality and Condorcet on OLLB come from 4/15/19/19 targets per rank bin and are correspondingly uncertain (with 19 targets a share of 0.16 has a 95\% Wilson interval of about 0.05--0.38).

\begin{table}[t]
\caption{Share of targets for which first place is infeasible, OLLB, cap \texttt{emp95}, by the target's mean-rule rank. Score rules: 4/15/80/4396 targets per bin; $^\ast$ rules run on the 57-target sample: 4/15/19/19.}
\label{tab:feas}
\centering\small
\begin{tabular}{@{}lrrrrr@{}}
\toprule
rule & 2--5 & 6--20 & 21--100 & $>$100 & smallest infeasible rank \\
\midrule
mean & 0 & 0 & 0 & 0.946 & 110 \\
Borda$^\ast$ & 0 & 0 & 0.158 & 0.579 & 22 \\
trimmed mean$^\ast$ & 0 & 0 & 0.263 & 1 & 37 \\
plurality$^\ast$ & 0 & 0.200 & 0.316 & 1 & 8 \\
median & 0 & 0.133 & 0.412 & 0.982 & 13 \\
geometric & 0 & 0.133 & 0.575 & 0.986 & 11 \\
harmonic & 0.250 & 0.533 & 0.662 & 0.991 & 5 \\
min & 0.250 & 0.733 & 0.775 & 0.991 & 5 \\
Condorcet$^\ast$ & 0 & 0.533 & 0.789 & 1 & 6 \\
\bottomrule
\end{tabular}
\end{table}

\begin{figure}[t]
\centering
\includegraphics[width=0.9\linewidth]{../figures/fig2_cap_sensitivity.pdf}
\caption{Cap sensitivity for the 20 closest challengers (mean-rule rank 2--21), the only set on which every rule was run on every target. Top: fraction feasible (95\% Wilson). Bottom: median $K$ over feasible targets (95\% bootstrap CI); its composition changes as feasibility drops. Cap families (ceiling, rival front, fixed gain, empirical) are not ordered against each other. $^\dagger$HELM empirical caps are transferred from OLLB.}
\label{fig:caps}
\end{figure}

\subsection{Targeted versus uniform improvement}
\label{sec:exp-uniform}
The uniform plan raises every task by the same amount $u^\ast$; it is the min-max ($L_\infty$) plan, so $L_\infty=u^\ast$ in every row (maximum deviation $10^{-6}$).
The comparison is therefore the volume ratio $L_1(\text{targeted})/L_1(\text{uniform})\le 1$.
For the mean it is exactly 1 in every group, dataset and cap: any plan reaching the same total costs the same, targeting only concentrates the points (OLLB: $K=1$ against 6 tasks; HELM ceiling, rank 6--20: $K=8$ [7, 9] against 57).
For every other rule targeting saves points: OLLB, ceiling, rank 6--20: min 0.27 [0.21, 0.30], plurality 0.25 [0.22, 0.30], median 0.33, Borda 0.36 [0.34, 0.39], Condorcet 0.37 [0.33, 0.41], harmonic 0.44 [0.43, 0.45], geometric 0.60 [0.58, 0.64], trimmed mean 0.67; HELM: min 0.044, median 0.111, plurality 0.111, Condorcet 0.32, Borda 0.60, harmonic 0.76, geometric 0.87 [0.86, 0.88] (Figure~\ref{fig:uniform}, Appendix~\ref{app:more}).
Targeting is close to uniform (more than half of targets at ratio $\ge 0.9$) only for the geometric mean under tight caps (OLLB \texttt{emp95} 0.95 [0.89, 0.97]; HELM $+10$ 0.93 [0.91, 0.96]) and, on HELM $+10$, the harmonic mean (0.87, 40\% of targets).
Rank 2--5 groups have three or four targets and wide intervals.


\subsection{Robustness and conditions of validity}
\label{sec:exp-robust}
We summarise the stability of the rule ordering by Kendall's $\tau$ between the order of rules by median $L_1$ (20 closest challengers) in a variant and in the main run (Figure~\ref{fig:robust}a,b).
\textbf{Normalisation.} On raw OLLB scores the ordering changes substantially ($\tau=0.21$ under the ceiling, 0.43 under a $+10$ cap): the median rule moves from cheapest (3.6) to 16.5 [14.6, 17.9], the mean from 29.2 to 19.0 [15.4, 19.7], and the arithmetic mean is no longer the most expensive rule; Borda stays among the cheapest.
\textbf{Leaderboard slices.} Removing merged models ($\tau=1.0$ under the ceiling, 0.79 under $+10$) or keeping the top 500 (0.86, 0.64) changes head-20 medians by at most 2.8 points (no merges) or, in the top-500 slice, by up to 7.4 points (plurality 9.0 against 16.4), and swaps are between rules whose intervals overlap; the ordering of Section~\ref{sec:exp-rules} holds on these slices.
\textbf{Leaving one task out} (OLLB, challengers redefined on the remaining five tasks): $\tau$ against the full six-task ordering, ceiling, is 0.93 (BBH), 0.79 (IFEval), 0.71 (MATH Lvl~5), 0.57 (GPQA), 0.43 (MMLU-Pro) and $-0.07$ for MuSR; dropping MuSR reorders the rules almost completely, so the six-task ordering depends on which task the challengers are strong on.
\textbf{Number of tasks.} On HELM, in random task subsets of size $n\in\{3,6,12,24,48\}$ (50 draws each, $+10$ cap, targets of rank 5 and 2 within each subset), the median $L_1$ of the rank-5 challenger grows roughly linearly in $n$ for the mean (log-log slope 1.06) and the geometric mean (1.02), with $K$ slopes 0.94 and 0.92, while for median and min $L_1$ is flat in $n$ ($-0.07$, $-0.08$) and $K$ stays between 1 and 2.5; Borda, Condorcet and plurality have $L_1$ slopes 0.89, 0.62 and $-0.17$. Slopes for the rank-2 challenger are noisy (some negative) and we make no claim of a scaling law from two ranks and no CI on these slopes.
\textbf{Noise.} With evaluation noise (binomial standard error per cell, 100 draws, ten closest challengers, $+10$ cap) the coefficient of variation of $L_1$ is 0.17 [0.13, 0.53] (HELM) and 0.20 [0.18, 0.48] (OLLB) for the mean but 0.68 and 0.65 for the median and 0.75 (44\% of draws infeasible) and 0.20 (51\% infeasible) for min; under a task bootstrap on HELM it is 0.22 for the mean and 0.92 for min.
The rules with the smallest $L_1$ (median, min) are the least stable, so their advantage in Section~\ref{sec:exp-rules} is smaller than the sampling variation of a single target's cost.


\begin{figure}[tb]
\centering
\includegraphics[width=\linewidth]{../figures/fig4_robustness_compact.pdf}
\caption{Robustness. (a) Median $L_1$ (95\% CI) for the 20 closest OLLB challengers under the ceiling, for normalised scores (main), raw scores, without merged models and in the top-500 slice. (b) Kendall $\tau$ of the rule order with one task left out against all six tasks. (c) HELM: median $K$ against the number of tasks in random subsets ($+10$ cap; solid rank-5, dashed rank-2 target). (d) Coefficient of variation of $L_1$ over draws under evaluation noise and task bootstrap (ten closest challengers); red column: percentage of infeasible draws for the numbered series (1), (2), (3) of the legend.}
\label{fig:robust}
\end{figure}

\subsection{Further checks}
\label{sec:exp-checks}
Details are in Appendix~\ref{app:checks}. The Condorcet-winner bound is tight within $\varepsilon$ against exact Copeland on HELM (not tested on OLLB). Under a convex (quadratic) cost the optimal plan touches all tasks, so ``concentrate on one task'' is an artefact of a linear cost. Random weights change the mean's cost by a factor of about two, and a tie margin of 0.1 instead of $10^{-3}$ changes median $L_1$ by at most 3.6\% but can shift HELM $K$ by a task or two.
